% Part: turing-machines 
% Chapter: machines-computations 
% Section: introduction

\documentclass[../../../include/open-logic-section]{subfiles}

\begin{document}

\olfileid{tur}{mac}{int}
\olsection{مقدمه}

محاسبه‌پذیر بودن تابعی، مثلاً از $\Nat$ به $\Nat$، به چه معناست؟ یکی
از نخستین و مشهورترین پاسخ‌ها این است که تابع هنگامی محاسبه‌پذیر است
که ماشین تورینگی بتواند آن را محاسبه کند. آلن تورینگ این مفهوم را در
سال ۱۹۳۶ صورت‌بندی کرد. ماشین‌های تورینگ نمونه‌ای از \emph{مدل
محاسبه} هستند---راهی به‌لحاظ ریاضی دقیق برای تعریف ایدهٔ «روش
محاسباتی». دربارهٔ معنای دقیق آن بحث هست، اما به‌طور گسترده پذیرفته
شده است که ماشین‌های تورینگ یکی از راه‌های مشخص‌کردن روش‌های محاسباتی
هستند. گرچه اصطلاح «ماشین تورینگ» تصویر دستگاهی فیزیکی با اجزای
متحرک را تداعی می‌کند، به‌معنای دقیق کلمه ماشین تورینگ سازه‌ای کاملاً
ریاضی است و ازاین‌رو ایدهٔ روش محاسباتی را آرمانی می‌کند. برای نمونه،
هیچ محدودیتی بر زمان یا حافظهٔ لازم برای ماشین تورینگ نمی‌گذاریم:
ماشین تورینگ می‌تواند چیزی را محاسبه کند حتی اگر محاسبه بیش از شمار
اتم‌های جهان فضای ذخیره‌سازی یا گام لازم داشته باشد.

\begin{explain}
شاید بهتر باشد ماشین تورینگ را برنامه‌ای برای نوع خاصی از سازوکار
خیالی بدانیم. این سازوکار از یک \emph{نوار} و یک \emph{سر
خواندن--نوشتن} تشکیل شده است. در نسخهٔ ما از ماشین تورینگ، نوار در
یک جهت (به راست) نامتناهی است و به \emph{خانه‌هایی} تقسیم می‌شود که
هریک ممکن است نمادی از \emph{الفبایی} متناهی داشته باشد. چنین
الفباهایی می‌توانند هر تعداد نماد متفاوت داشته باشند، اما ما عمدتاً
به سه نماد بسنده می‌کنیم: $\TMendtape$، $\TMblank$ و~$\TMstroke$.
هنگام آغاز سازوکار، نوار تهی است (یعنی هر خانه نماد $\TMblank$ را
دارد)، جز چپ‌ترین خانه که $\TMendtape$ را دارد و شمار متناهی‌ای از
خانه‌ها که \emph{ورودی} را دربر می‌گیرند. سازوکار در هر لحظه در یکی
از شمار متناهی \emph{حالت‌ها} است. در آغاز، سر چپ‌ترین خانه را در
\emph{حالت آغازین} مشخصی می‌خواند. در هر گام اجرای سازوکار، محتوای
خانه‌ای که اکنون خوانده می‌شود، همراه با حالت سازوکار و برنامهٔ ماشین
تورینگ، رخداد بعدی را تعیین می‌کنند. برنامهٔ ماشین تورینگ با تابعی
جزئی داده می‌شود که حالت~$q$ و نماد~$\sigma$ را ورودی می‌گیرد و سه‌تایی
$\tuple{q',\sigma',D}$ را بیرون می‌دهد. هرگاه سازوکار در حالت $q$
باشد و نماد $\sigma$ را بخواند، نماد خانهٔ کنونی را با $\sigma'$
جایگزین می‌کند؛ سر بسته به اینکه $D$ برابر $\TMleft$، $\TMright$ یا
$\TMstay$ باشد به چپ، راست یا هیچ‌جا حرکت می‌کند؛ و سازوکار به حالت
$q'$ می‌رود.

برای نمونه، وضعیت \olref{fig:tm} را در نظر بگیرید.
\begin{figure}
  \olasset{\olpath/assets/diagrams/turing-machine.tikz}
  \caption{ماشین تورینگی در حال اجرای برنامهٔ خود.}
  \ollabel{fig:tm}
\end{figure}
بخش نمایان نوار ماشین تورینگ در چپ‌ترین خانه نماد پایان نوار
$\TMendtape$ دارد و پس از آن سه $1$، یک $0$ و چهار $1$ دیگر آمده است.
سر خانهٔ سوم از چپ را می‌خواند که یک~$1$ دارد و در حالت~$q_1$ است---
می‌گوییم «ماشین در حالت~$q_1$ یک $1$ می‌خواند.» اگر برنامهٔ ماشین
تورینگ برای ورودی $\tuple{q_1,1}$ سه‌تایی
$\tuple{q_2,0,\TMstay}$ را بازگرداند، ماشین اکنون~$1$ خانهٔ سوم را
با~$0$ جایگزین می‌کند، سر خواندن--نوشتن را همان‌جا می‌گذارد و به
حالت~$q_2$ می‌رود. اگر سپس برنامه برای ورودی $\tuple{q_2,0}$ سه‌تایی
$\tuple{q_3,0,\TMright}$ را بازگرداند، ماشین~$0$ را با~$0$ دیگری
بازنویسی می‌کند (در عمل محتوای نوار زیر سر خواندن--نوشتن را تغییر
نمی‌دهد)، یک خانه به راست می‌رود و وارد حالت~$q_3$ می‌شود؛ و به همین
ترتیب ادامه می‌دهد.

می‌گوییم ماشین هنگامی \emph{متوقف می‌شود} که به حالت~$q_n$ و نماد
$\sigma$ای برسد که هیچ دستوری برای $\tuple{q_n,\sigma}$ وجود ندارد؛
یعنی تابع گذار روی ورودی $\tuple{q_n,\sigma}$ تعریف نشده است. به بیان
دیگر، ماشین دستوری برای اجرا ندارد و در آن نقطه کار را پایان می‌دهد.
توقف را گاه با حالت توقف خاصی چون~$h$ نشان می‌دهند. بعداً این را با
جزئیات بیشتر خواهیم دید.
\end{explain}

\begin{digress}
زیبایی مقالهٔ تورینگ، «دربارهٔ اعداد محاسبه‌پذیر»، در این است که او
نه‌تنها تعریفی صوری، بلکه استدلالی ارائه می‌کند که نشان می‌دهد تعریف
مفهوم شهودی محاسبه‌پذیری را به‌درستی می‌گیرد. از تعریف باید روشن باشد
هر تابعی که ماشین تورینگ محاسبه می‌کند به‌معنای شهودی محاسبه‌پذیر
است. تورینگ سه نوع استدلال برای درستی عکس آن عرضه می‌کند؛ یعنی اینکه
هر تابعی که طبیعتاً محاسبه‌پذیر می‌دانیم با چنین ماشینی محاسبه‌پذیر
است. این‌ها (به تعبیر تورینگ) عبارت‌اند از:
\begin{enumerate}
\item استناد مستقیم به شهود.
\item اثبات هم‌ارزی دو تعریف (اگر تعریف تازه جاذبهٔ شهودی بیشتری داشته باشد).
\item ارائهٔ نمونه‌هایی از رده‌های بزرگ اعداد محاسبه‌پذیر.
\end{enumerate}
هدف ما تعریف مفهوم محاسبه‌پذیری «در اصل» است؛ یعنی بدون درنظرگرفتن
محدودیت‌های عملی زمان و فضا. البته پس از تثبیت گسترده‌ترین تعریف
محاسبه‌پذیری می‌توان محاسبه با منابع کراندار را بررسی کرد؛ این هستهٔ
موضوعی است که «پیچیدگی محاسباتی» نام دارد.
\end{digress}

\begin{history}
آلن تورینگ ماشین‌های تورینگ را در ۱۹۳۶ ابداع کرد. گرچه آن زمان به
تصمیم‌پذیری منطق مرتبهٔ اول علاقه داشت، مقاله‌اش مقاله‌ای تعیین‌کننده
دربارهٔ مبانی طراحی رایانه وصف شده است. تورینگ در مقاله بر اعداد حقیقی
محاسبه‌پذیر، یعنی اعداد حقیقی با بسط اعشاری محاسبه‌پذیر، تمرکز می‌کند؛
اما یادآور می‌شود سازگارکردن مفاهیمش با توابع محاسبه‌پذیر روی اعداد
طبیعی و جز آن دشوار نیست. توجه کنید این پنج سال کامل پیش از ساخت
نخستین رایانهٔ همه‌منظورهٔ عملی در ۱۹۴۱ بود (به دست کنراد تسوزهٔ
آلمانی در اتاق نشیمن خانهٔ والدینش)، هفت سال پیش از ساخت رایانهٔ
رمزشکن کلوسوس به دست تورینگ و همکارانش در بلچلی پارک (۱۹۴۳)، نه سال
پیش از انیاک آمریکایی (۱۹۴۵)، دوازده سال پیش از ساخت نخستین رایانهٔ
همه‌منظورهٔ بریتانیایی---ماشین آزمایشی مقیاس کوچک منچستر---در منچستر
(۱۹۴۸)، و سیزده سال پیش از نخستین آزمایش بیناک به دست آمریکایی‌ها
(۱۹۴۹). ماشین آزمایشی منچستر نخستین رایانهٔ دارای برنامهٔ ذخیره‌شده
بود---ماشین‌های پیشین برای هر کار تازه باید دستی سیم‌کشی مجدد می‌شدند.
\end{history}

\end{document}
